\section{Bài tập 9: AC/DC Power Circuit Application}

Nguồn xoay chiều $V_1$: \texttt{VOFF} = $0\text{ V}$, \texttt{VAMPL} = $220\text{ V}$, \texttt{FREQ} = $50\text{ Hz}$, \texttt{AC} = $1$; biến áp $TX_1$ tỷ số $10:1$, cầu chỉnh lưu 4 diode \texttt{Dbreak}, tải $R_1 = 1\text{ k}\Omega$. Mô phỏng Transient trong $100\text{ ms}$, \textit{Maximum step size} $= 0.1\text{ ms}$.

\Needspace{0.4\textheight}
\subsection{Bước 1: Điện áp chỉnh lưu chưa lọc, chưa ổn áp}
\given{ex9_given}{0.65\textwidth}{Sơ đồ Bước 1: mạch chỉnh lưu chưa có mạch lọc (Hình 2.29)}

\simresult
\simfig{ex9_step1_sim}{Kết quả mô phỏng dạng sóng điện áp ngõ ra tại Bước 1}

\begin{notebox}{Nhận xét, Bước 1}
  \begin{itemize}[leftmargin=1.5em,itemsep=2pt,topsep=2pt]
    \item Số liệu đọc bằng con trỏ: $V_{\max} = 21.105\text{ V}$ (tại $t = 15\text{ ms}$), $V_{\min} \approx 0.156\text{ V} \approx 0\text{ V}$ (tại $t \approx 30\text{ ms}$), $V_r \approx 20.95\text{ V}$.
    \item Điện áp trên tải là điện áp một chiều nhấp nhô, các đỉnh cách nhau $10\text{ ms}$ (tần số gợn sóng $100\text{ Hz}$) và luôn không âm.
    \item Cuối mỗi nửa chu kỳ điện áp về $0\text{ V}$, nên dạng sóng này chưa dùng trực tiếp để cấp nguồn cho mạch điện tử.
  \end{itemize}
\end{notebox}

\Needspace{0.4\textheight}
\subsection{Bước 2: Thêm tụ lọc \texorpdfstring{$C_1 = 10\,\mu\text{F}$}{C1 = 10uF}}
\given{ex9_step2_given}{0.65\textwidth}{Sơ đồ Bước 2: thêm tụ lọc $C_1 = 10\,\mu\text{F}$ (Hình 2.30)}

\simresult
\simfig{ex9_step2_sim}{Kết quả mô phỏng dạng sóng ngõ ra tại Bước 2 ($C_1 = 10\,\mu\text{F}$)}

\begin{notebox}{Nhận xét, Bước 2}
  \begin{itemize}[leftmargin=1.5em,itemsep=2pt,topsep=2pt]
    \item Số liệu đọc bằng con trỏ: $V_{\max} = 21.110\text{ V}$ (tại $t = 15\text{ ms}$), $V_{\min} = 10.836\text{ V}$ (tại $t \approx 31.5\text{ ms}$), $V_r \approx 10.27\text{ V}$.
    \item Điện áp ngõ ra không còn về $0\text{ V}$ mà có dạng răng cưa: tụ nạp theo đỉnh chỉnh lưu, rồi xả dần qua $R_1$ khi các diode khóa.
    \item Hằng số thời gian $\tau = R_1 C_1 = 10\text{ ms}$ ngang với nửa chu kỳ ($10\text{ ms}$) nên tụ xả nhiều và gợn sóng còn lớn. Theo mô hình xả mũ $v_C(t) = V_p\,e^{-(t-t_0)/\tau}$, tụ xả khoảng $8\text{ ms}$ giữa hai đỉnh nên $V_{\min} \approx 10.8\text{ V}$ và $V_r \approx 20.6 - 10.8 \approx 9.8\text{ V}$.
  \end{itemize}
\end{notebox}

\Needspace{0.62\textheight}
\subsection{Bước 3: Đổi tụ lọc thành \texorpdfstring{$C_1 = 680\,\mu\text{F}$}{C1 = 680uF}}
\simresult
\simfig{ex9_step3_sim}{Kết quả mô phỏng dạng sóng ngõ ra tại Bước 3 ($C_1 = 680\,\mu\text{F}$)}

\begin{notebox}{Nhận xét, Bước 3}
  \begin{itemize}[leftmargin=1.5em,itemsep=2pt,topsep=2pt]
    \item Số liệu đọc bằng con trỏ (đặt tại $t = 75\text{ ms}$ và $t = 85\text{ ms}$): $V_{\max} = 20.684\text{ V}$, $V_{\min} = 20.420\text{ V}$, $V_r \approx 0.264\text{ V}$.
    \item Hằng số thời gian tăng lên $\tau = R_1 C_1 = 680\text{ ms} \gg 10\text{ ms}$, tụ xả rất chậm giữa hai đỉnh nên gợn sóng giảm mạnh (từ $10.27\text{ V}$ xuống khoảng $0.26\text{ V}$).
    \item Ước lượng gợn sóng: $V_r \approx \dfrac{I_{dc}}{2fC_1} \approx \dfrac{20\text{ mA}}{100 \times 680 \times 10^{-6}} \approx 0.29\text{ V}$.
    \item Điện áp ngõ ra gần như phẳng ở mức $\approx 20.5\text{ V}$ sau khoảng $15\text{ ms}$ đầu (giai đoạn tụ nạp).
  \end{itemize}
\end{notebox}

\Needspace{0.4\textheight}
\subsection{Bước 4: Thêm Diode Zener \texorpdfstring{$V_Z = 22\text{ V}$}{VZ = 22V}}
\given{ex9_step4_given}{0.65\textwidth}{Sơ đồ Bước 4: mạch chỉnh lưu có tụ lọc và Diode Zener (Hình 2.31)}

\simresult
\simfig{ex9_step4_sim}{Kết quả mô phỏng điện áp ngõ ra và dòng Zener tại Bước 4 ($V_Z = 22\text{ V}$)}

\begin{notebox}{Nhận xét, Bước 4}
  \begin{itemize}[leftmargin=1.5em,itemsep=2pt,topsep=2pt]
    \item Số liệu đọc bằng con trỏ ($t \approx 94\text{--}95\text{ ms}$): $V_{\max} = 20.652\text{ V}$, $V_{\min} = 20.146\text{ V}$, $V_r \approx 0.506\text{ V}$; dòng Zener $I_Z \approx 0.12\text{ mA}$ ($121$ đến $124\,\mu\text{A}$).
    \item Điện áp trên tụ (không quá $20.7\text{ V}$) nhỏ hơn $V_Z = 22\text{ V}$ nên Zener chưa bị đánh thủng, $I_Z \approx 0$.
    \item Zener chưa tham gia điều chỉnh, dạng sóng ngõ ra giống Bước 3.
  \end{itemize}
\end{notebox}

\Needspace{0.62\textheight}
\subsection{Bước 5: Đổi Diode Zener thành \texorpdfstring{$V_Z = 20\text{ V}$}{VZ = 20V}}
\simresult
\simfig{ex9_step5_sim}{Kết quả mô phỏng điện áp ngõ ra và dòng Zener tại Bước 5 ($V_Z = 20\text{ V}$)}

\begin{notebox}{Nhận xét, Bước 5}
  \begin{itemize}[leftmargin=1.5em,itemsep=2pt,topsep=2pt]
    \item Số liệu đọc bằng con trỏ ($t \approx 94\text{--}95\text{ ms}$): $V_{\max} = 20.081\text{ V}$, $V_{\min} = 19.746\text{ V}$, $V_r \approx 0.334\text{ V}$; dòng Zener đạt đỉnh $I_Z \approx 2.21\text{ A}$ tại $t = 94.9\text{ ms}$ và khoảng $0.13\text{ mA}$ ngoài xung.
    \item Điện áp đỉnh của tụ ($\approx 20.6\text{ V}$) lớn hơn $V_Z = 20\text{ V}$ nên Zener bị đánh thủng ở các đỉnh và ghim điện áp ngõ ra quanh $20\text{ V}$.
    \item Dòng Zener xuất hiện thành các xung ngắn tại $t \approx 15, 35, 55, 75, 95\text{ ms}$, đúng lúc điện áp tụ chạm đỉnh. Nhờ đó gợn sóng giảm từ $0.51\text{ V}$ (Bước 4) xuống $0.33\text{ V}$ và điện áp ngõ ra thấp hơn Bước 4 khoảng $0.5\text{ V}$.
  \end{itemize}
\end{notebox}

\subsection{Bảng tổng hợp kết quả}
\begin{table}[H]
  \centering
  \caption{Tổng hợp điện áp ngõ ra qua 5 bước (lý thuyết ước tính theo mô hình diode $V_F = 0.7\text{ V}$)}
  \begin{tabular}{lcccccc}
    \toprule
    & \multicolumn{3}{c}{\textbf{Lý thuyết (ước tính)}} & \multicolumn{3}{c}{\textbf{PSpice}} \\
    \cmidrule(lr){2-4} \cmidrule(lr){5-7}
    \textbf{Bước} & $V_{\max}\text{ (V)}$ & $V_{\min}\text{ (V)}$ & $V_r\text{ (V)}$ & $V_{\max}\text{ (V)}$ & $V_{\min}\text{ (V)}$ & $V_r\text{ (V)}$ \\
    \midrule
    1: không lọc               & $20.6$ & $0$    & $20.6$   & $21.105$ & $0.156$  & $20.949$ \\
    2: $C_1 = 10\,\mu\text{F}$ & $20.6$ & $10.8$ & $9.8$    & $21.110$ & $10.836$ & $10.273$ \\
    3: $C_1 = 680\,\mu\text{F}$& $20.6$ & $20.3$ & $0.3$    & $20.684$ & $20.420$ & $0.264$ \\
    4: Zener $22\text{ V}$     & $20.6$ & $20.3$ & $0.3$    & $20.652$ & $20.146$ & $0.506$ \\
    5: Zener $20\text{ V}$     & $\approx 20$ & $\approx 20$ & $\approx 0$ & $20.081$ & $19.746$ & $0.334$ \\
    \bottomrule
  \end{tabular}
\end{table}
